\documentclass[11pt, aspectratio = 169]{beamer}

\usetheme[progressbar=frametitle]{metropolis}
\usepackage{appendixnumberbeamer}
\usepackage{booktabs}
\usepackage[scale=2]{ccicons}
\usepackage{pgfplots}
\usepgfplotslibrary{dateplot}
\usepackage{xspace}
\usepackage{hyperref}
\usepackage{verbatim} 
\usepackage{comment}
\usepackage{cancel}
\definecolor{links}{HTML}{1B1B81}
\hypersetup{colorlinks,linkcolor=blue,urlcolor=blue}

% listings: To highlight code:
\usepackage{listings}
\lstset{
  language=R,
  basicstyle=\ttfamily\footnotesize,
  keywordstyle=\color{blue},
  commentstyle=\color{green!50!black},
  stringstyle=\color{red},
  frame=single,
  showstringspaces=false
}
\lstset{basicstyle=\ttfamily\footnotesize, frame=single}

\title{Applied Economics}
\subtitle{Causality and Correlation}
\author{Pello Aspuru}
\date{}
\institute{Winter 2025/26}
% \titlegraphic{\hfill\includegraphics[height=1.5cm]{logo.pdf}}
\setbeamercolor{normal text}{fg=black}
\setbeamercolor{block title}{fg=teal}
\setbeamercovered{transparent} % default faintness (~15%)

\let\svdescription\description
\let\svenddescription\enddescription
\renewenvironment{description}{%
  \svdescription%
  \let\svitem\item%
  \renewcommand\item[1][]{\svitem[\textcolor{\descriptioncolor}{##1}]}%
  }{
  \svenddescription%
  }
\newcommand\descriptioncolor{teal}

%\usepackage{Sweave}
\begin{document}
%\Sconcordance{concordance:AE_w7_CC.tex:AE_w7_CC.Rtex:%
1 36 1 1 0 395 1}


\maketitle

\begin{frame}{Plan for this Lecture}
\begin{enumerate}

\item Correlation vs. Causality
\item Endogeneity
\begin{itemize}
    \item Reverse Causality
    \item Measurement Error
    \item Omitted Variable Bias
\end{itemize}
\item Model Specification

\end{enumerate}
\end{frame}

\section{Correlation vs. Causality}

\begin{frame}{Correlation vs Causation}
\linespread{1.3}\selectfont % dense slide: tighter line spacing so it fits

\textcolor{teal}{Correlation} is a \textbf{statistical measure} describing the \textbf{relation} between treatment variable $x$ and outcome variable $y$.\\
    \begin{itemize}
    \item[] \textit{"Those who have a university degree earn more that those who do not"}
    \item[] \textit{"When icecream consumption increases, sea drownings increase"}
    \item[] \textit{"Healthier people are more wealthy"}
\end{itemize}
\onslide<2->{
\textcolor{teal}{Causation} is a \textbf{claim} affirming that treatment $x$ causes outcome $y$.
\begin{itemize}
    \item[] \textit{"Having a university degree makes you earn more"} \onslide<5->{Selection Bias}
    \item[] \textit{"Icecream consumption makes pepole drown in the sea"} \onslide<6->{Ommitted Variable}
    \item[] \textit{"Better health makes people more wealthy"} \onslide<7->{Reverse Causality}
\end{itemize}}
%Causality: $x$ \textbf{causes} $y$  $\Leftrightarrow$ if we were to intervene and change the value of variable $x$ \textit{without changing anything else}, 
%then $y$ would also change as a result.

\onslide<3->{
A correlation \textbf{does not necessarily} reflect a causal relationship. Why?}
\begin{itemize}
    \onslide<4->{
    \item[]  Endogeneity: Selection bias, omitted variables, reverse causality\dots}
\end{itemize}
%(check \href{https://tylervigen.com/spurious-correlations/}{Spurious Correlations})

\end{frame}

\begin{frame}{Causal Inference}

Causal inference consists on using econometric tools to \textbf{credibly} support causal claims.

\begin{enumerate}
    \onslide<2-4>{
    \item First-best: field experiments or \textbf{randomised control trials (RCTs)} 
    \begin{itemize}
        \item A ``treatment'' is randomly allocated to find its causal effect on an outcome.
        \item Compare mean outcome of treatment and control groups to infer causal effect.
        %\item Example: random allocation of deworming treatment on Kenyan children to study school attendance. Miguel, E., & Kremer, M. (2004).
        \item Popular in development economics, but often unethical or impractical in other fields.
    \end{itemize}}
    \onslide<3->{
    \item Second-best: quasi-experiments or natural experiments.
    \begin{itemize}
        \item Setups where external events, policies, or random conditions created a random allocation of the treatment.
        %\item Example: compare fast-food restaurants along a US-state frontier, as one state had higher minimum wage than the other. Card, David & Krueger, Alan B, (1994)  
    \end{itemize}}
    \onslide<4->{
    \item Third best: non-experiments, observational data.
    \begin{itemize}
        \item Non-random allocation of treatment. Claiming causality becomes tricky.
        \item Use multiple linear regression (MLR) to reduce endogeneity concerns.
    \end{itemize}}
\end{enumerate}

\onslide<5->{This course will focus on MLR and quasi-experimental methods.}

\end{frame}

\begin{frame}[fragile]{Big Picture: Steps in empirical economic analysis}
    \begin{center}
        Economic question: does education \alert{cause} an increase in wages?
    \end{center}
    \vspace{-3mm}
\pause
\vspace{+0.2cm}
\begin{enumerate}
    \item Construct an economic model for the \alert{representative agent}.
        $$wage = \beta_0 + \beta_1 educ + \beta_2 expert + \beta_3 ability $$
    \vspace{-4mm}
    \item Turn the economic model into a statistical model for the \alert{population}.
        $$wage_{i} = \beta_0 + \beta_1 educ_{i} + \beta_2 expert_{i} + \beta_3 ability_{i} + \alert{u_{i}} \quad i=1,...,N$$
    \pause
    \vspace{-4mm}
    \item Turn the statistical model for the \alert{observed sample}.
        $$wage_{i} = \hat{\beta}_0 + \hat{\beta}_1 educ_{i} + \hat{\beta}_2 expert_{i} + \hat{\beta}_3 ability_{i} + \alert{\hat{u}_{i}} \quad i=1,...,n$$
    \pause
    \vspace{-4mm}
    \begin{lstlisting}[language=R]
    df_wages <- read.csv("./AE_week7_CC/wage2.csv")
    mlr <- lm(wage ~ educ + exper +  IQ, data = df_wages)
    summary(mlr)
    \end{lstlisting}
\end{enumerate}
\end{frame}

\begin{frame}{Causality}
We say $X$ \alert{causes} $Y$ if: 

When we change the value of $X$ \textbf{without changing anything else}, then $Y$ would also change as a result.
\pause
\\
An ideal way to identify causal effects are \color{teal}\textbf{Randomized Control Trials} (RCTs)
\begin{itemize}
    \item you randomly assign people to a treatment and a control group. \\
  $\rightarrow$ the treatment assignment is independent of all other potential factors that \hspace*{+0.32cm} have an impact on $Y$. \\
$\rightarrow$ treatment effects can be estimated by simply comparing means across groups
\pause
\item example?
\end{itemize}
\end{frame}


\begin{frame}{Causal Analysis with Observational Data}

RCT disadvantages
\begin{itemize}
    \item Expensive.
    \item Ethical issues.
    \item External validity issues: to what extent results from a (small) RCT can be generalized to other contexts (e.g. other countries).
\end{itemize}
\pause
Thus, we often use observational data (as opposed to experimental data).
\\
With this type of data, causal inference is usually more difficult.

Coefficients in a (multiple) linear regressions measure \alert{(partial) correlations} between variables.
\alert{They can, but they don’t have to reflect causal effects!}
\end{frame}

\begin{comment}

\begin{frame}{Correlation vs. Causality: Example}

Regress height of son ($S_i$) on the height of his father ($F_i$), after standardizing both variables ($\rightarrow$ mean = 0, sd = var = 1): $S_i = \beta_0 + \beta_1F_i + u_i$
\vspace{-0.15cm}
\begin{itemize}
    \item $\beta_1$ might plausibly reflect a causal effect of paternal height on height of the son. 
    \item $\hat{\beta}_1 = \frac{\widehat{cov}(S_i, F_i)}{\widehat{var}(F_i)} = \widehat{cov}(S_i, F_i) = \widehat{cor}(S_i, F_i) $
\end{itemize}
\pause
We can flip dependent and independent variable: $F_i = \delta_0 + \delta_1S_i + e_i$
\vspace{-0.15cm}

\begin{itemize}
    \item $\delta_1$ cannot plausibly reflect a causal effect
    \item $\hat{\delta}_1 = \frac{\widehat{cov}(F_i, S_i)}{\widehat{var}(S_i)} = \widehat{cov}(F_i, S_i) = \widehat{cor}(F_i, S_i) = \hat{\beta_1}$
\end{itemize}

\end{frame}


\begin{frame}{Correlation vs. Causation: Examples}
\color{teal}\textbf{Take-away:} 


\color{black} Neither estimated nor true slope parameters necessarily reflect a causal effect. 

Regressions per se are just tools to analyse (partial)  correlations between variables.


Causality is an assumption that we add to it, and that needs to be based on economic theory/common sense.
\end{frame}


\begin{frame}{Correlation vs. Causation: More examples}
Even if we have a causal model in the background, there are several reasons why correlations in observational data might not reflect the causal effect of interest:
\begin{itemize}
    \item the observed correlation between  \alert{\textit{education}} and the \alert{\textit{wage}} might not reflect the causal effect of education, because individuals who choose to get more education might differ from those who don't (higher innate ability, motivation) \alert{[omitted variables]}
    \item the observed correlation between the \alert{\textit{minimum wage level}} and the \alert{\textit{employment rate}} might not reflect the causal effect of the minimum wage because policymakers may increase the minimum wage in times when the employment rate is high \alert{[reverse causality/simultaneity]}
\end{itemize}
\end{frame}
\end{comment}

\section{Exogeneity and Endogeneity}

\begin{frame}{Endogeneity}
Exogeneity (Assumption 4) is central to obtain causal effects.
$$E(u_i|\boldsymbol{X}_i) = E(u_i) = 0$$
\begin{itemize}
    \item Regressors $\boldsymbol{X}_i$ and the error term $u_i$ are independent.
    \item The regressors $\boldsymbol{X}_i$ behave as if they had been randomly assigned. 
\end{itemize}
\pause
\begin{block}{Endogeneity $E(u_i|\boldsymbol{X}_i) \neq 0$}
    \vspace{2mm}
    Correlation between one or more regressors $\boldsymbol{X}_i$ and the error term $u_i$. \\
    Regressors $\boldsymbol{X}_i$ are \alert{not} as if randomly assigned.
\end{block}
\pause
Implication: OLS estimator is biased and inconsistent: 
\begin{align*}
E[\hat{\boldsymbol{\beta} }] = \boldsymbol{\beta}  + E[(\boldsymbol{X}'\boldsymbol{X})^{-1}(\boldsymbol{X}'\boldsymbol{u})]  \neq \boldsymbol{\beta}
\end{align*}
\end{frame}

\begin{frame}[fragile]{Endogeneity Example in a Generic Simple Regression Model}  
     \vspace{-4mm}  
    \begin{align*}
        Y_{it} = \beta_0 + \beta_1 X_{it} + \alert{u_{it}} \quad i=1,...,N
    \end{align*}
    \vspace{2mm}
    \pause
    Parameter estimate for the observed sample is:
    \begin{align*}
        \hat{\beta}_1 = \beta_1 + \frac{\frac{1}{n} \sum_{i=1}^n(X_i - \overline{X}) \cdot u_i}{\frac{1}{n} \sum_{i=1}^n(X_i - \overline{X})^2}
    \end{align*}
    \vspace{1mm}
\pause
Under Assumptions 1 to 3:
\[
\left.
\begin{array}{l}
\bullet\ \frac{1}{n} \sum_{i=1}^n (X_i - \overline{X}) \cdot u_i 
    \xrightarrow{p} \operatorname{cov}(X_i, u_i) \\
\bullet\ \frac{1}{n} \sum_{i=1}^n (X_i - \overline{X})^2 
    \xrightarrow{p} \operatorname{var}(X_i)
\end{array}
\right\}
\qquad 
\hat{\beta}_1 \xrightarrow{p} \beta_1 + \frac{\alert{\operatorname{cov}(X,u)}}{\operatorname{var}(X)}
\]
\pause
$\Rightarrow$ \textbf{Direction of the inconsistency depends on the sign of \alert{$\operatorname{cov}(X,u)$}}
\end{frame}

\begin{frame}{Sources of Endogeneity}
Frequent reasons for endogeneity are
\begin{itemize}
    \item Omitted variables
    \item Measurement error in the explanatory variable
    \item Reverse causality.
\end{itemize}
\end{frame}

\section{Omitted Variables}

\begin{frame}[fragile]{Omitted Variables Example}
\vspace{-2mm}
\[
\text{Relevant model: }
\begin{aligned}
    \quad wage_i = \beta_0 + \beta_1\, educ_i + \beta_2\, ability_i + u_i
\end{aligned}
\]
\begin{lstlisting}[language=R]
    mlr <- lm(log(wage) ~ educ +  IQ, data = df_wages)
\end{lstlisting}
\pause
\[
\text{Estimated model: }
\begin{aligned}
        \quad wage_{i} = \beta_0 + \beta_1 educ_{i} + \cancel{\beta_2 ability_{i}} + v_i \quad v_i =  \alert{\beta_2} \cdot ability_{i} + u_i
\end{aligned}
\]
\begin{lstlisting}[language=R]
    slr <- lm(log(wage) ~ educ, data = df_wages)
\end{lstlisting}
\pause
Endogeneity requires:
\begin{itemize}
    \item $\alert{\beta_2} \neq 0$: omitted variable to correlate with outcome.
    \item $E[educ_i \cdot ability_i] \neq 0$: omitted variable to correlate with the regressor.
\end{itemize}
 \begin{align*}
    E[educ_{i} \cdot v_i] 
     %= E[educ_{i} \cdot \beta_2 \cdot ability_{i} + educ_{i} \cdot u_i]
     = \alert{\beta_2} \cdot E[educ_i \cdot ability_i] + \cancelto{0}{E[educ_i \cdot u_i]}
\end{align*}
\end{frame}

\begin{frame}[fragile]{Direction of the Bias}

Simmilarly, the covariance between error term and regresson can be expressed as:
\begin{align*}
    \operatorname{Cov}(X,u) \rightarrow \operatorname{Cov}(educ,v) = \beta_2 \operatorname{Cov}(educ,ability) + \cancelto{0}{\operatorname{Cov}(educ,u)}
\end{align*}
\begin{align*}
    \hat{\beta}_1 \xrightarrow{p} \beta_1 +
    \frac{\operatorname{Cov}(X,u)}{\operatorname{Var}(X)} \quad \equiv \quad \beta_1 + \beta_2 \cdot \frac{\operatorname{Cov}(educ,ability)}{\operatorname{Var}(educ)}
\end{align*}
\pause
\vspace{2mm}
Direction of the \textbf{Omitted Variable Bias} depends on:
\vspace{-3mm}
\begin{itemize}
    \item the sign of $\beta_2$ (impact of ability on wages).
    \begin{lstlisting}[language=R]
    test_1 <- lm(log(wage) ~ IQ, data = df_wages)
    \end{lstlisting}
    \item the direction of $\operatorname{Cov}(educ, ability)$ (the impact of ability on education)
    \begin{lstlisting}[language=R]
    test_2 <- lm(educ ~ IQ, data = df_wages)
    \end{lstlisting}
    \pause
    \vspace{1mm}
    \item[$\rightarrow$] Will the parameter estimate be upward or downward biased?
\end{itemize}

\end{frame}

\begin{frame}[fragile]{Omitted Variables: Example}
In the Esitmated Model, $\hat{\beta}_1$ is \textbf{\color{teal}biased upward} because 
\begin{itemize}
    \item $\beta_2 > 0$: ability increases wages.
    \item $cov(IQ_i, educ_i) > 0$: $education$ is positively correlated with the omitted variable
    \footnotesize

\end{itemize}
Thus, we \textit{overestimate} the true positive impact of education on wages; part of the correlation reflects the \textit{indirect} effect that individuals with higher education have a higher IQ, and therefore have higher wages.

\bigskip

In real-world applications, we do not have data on the omitted variable. \\
We can only speculate about its correlation with the included regressor(s) and its impact on the dependent variable, to think about how its omission affects the parameter of interest.
\end{frame}

\section{Measurement Error}

\begin{frame}{(Classical) Measurement Error in the Regressor}
Original model: $\quad \text{wage}_i = \beta_0 + \beta_1 \text{educ}_i + u_i  \quad E[u_i \mid \text{educ}_i] = 0$ 
\\
\pause
We do not observe $\text{educ}_i$ but just a noisy measure $\widetilde{\text{educ}}_{i}$:
$\quad \widetilde{\text{educ}}_{i} = \text{educ}_{i} + \alert{w_i}$
\vspace{-2mm}
\begin{itemize}
    \item $E[w_i] = 0$
    \item $E[w_i \text{educ}_i] = 0$ 
    \item $E[w_i u_i] = 0$
\end{itemize}
Estimated model: 
$\quad \text{wage}_i = \beta_0 + \beta_1 \widetilde{\text{educ}}_{i} + e_i \quad \quad e_i = - \beta_1 w_{i} + u_i$
\vspace{+0.2cm}
\pause
\begin{align*}
Cov(X, u)  = Cov(\widetilde{\text{educ}}_{i}, e_i) 
= & \ E(e_i \widetilde{\text{educ}}_{i}) \\
= & \ E[(-\beta_1 w_{i} + u_i)\widetilde{\text{educ}}_{i}] \\
= & \ -\beta_1 E[w_i \widetilde{\text{educ}}_{i}] + E[u_i \widetilde{\text{educ}}_{i}] \\
= & \ -\beta_1 E[w_i \text{educ}_i + w_i^2] + E[u_i \text{educ}_i + u_i w_i] \\
= & \ -\beta_1 E[w_i^2] \neq 0
\end{align*}
\end{frame}

\begin{frame}{Measurement Error causes attenuation bias}
\begin{align*}
\hat{\beta}_1
\quad \xrightarrow{p} \quad & \beta_1 + \frac{ Cov(\widetilde{\text{educ}}_{i}, e_i) }{ Var(\widetilde{\text{educ}}_{i}) } \\
&= \beta_1 - \beta_1 \frac{ E[w_i^2] }{ Var(\text{educ}_i + w_i) } \\
&= \beta_1 - \beta_1 \frac{ Var(w_i) }{ Var(\text{educ}_i) + Var(w_i) } \\
&= \beta_1 \frac{ Var(\text{educ}_i) + Var(w_i) - Var(w_i) }{ Var(\text{educ}_i) + Var(w_i) } \\
&= \beta_1 \frac{ Var(\text{educ}_i) }{ Var(\text{educ}_i) + Var(w_i) }
\end{align*}

$0 < \frac{ Var(\text{educ}_i) }{ Var(\text{educ}_i) + \alert{Var(w_i)}} \leq 1 
\Rightarrow \hat{\beta}_1$ biased towards zero (attenuation bias)
\end{frame}


\section{Reverse Causality}

\begin{frame}{Reverse Causality}
We talk about reverse causality when $Y$ causes $X$ \textbf{and} $X$ causes $Y$.
\pause
\textbf{Previous Example:} \\
\begin{itemize}
    \item  We are interested in the effect of school budget on student outcomes: $$grades_i = \beta_0 + \beta_1 funding_i + u_i$$

\item But the grades also determine government spending on a school: $$funding_i = \delta_0 + \delta_1 grades_i + v_i$$ 

\item Assume: $cov(u_i, v_i ) = 0$

\item Still, we get endogeneity: $cov(funding_i, u_i ) \neq 0$, \\
because $pol \,funding_i$ (linearly) depends on $grades_i $ and $grades_i$ depends on $u_i$ 
\end{itemize}
\end{frame}


\begin{comment}
\section{Model Specification}

\begin{frame}{How to choose which regressors to include in your model?}
If our aim is to model causal relationships, the choice of which regresssors to include should 
\begin{itemize}
    \item be based primarily on economic arguments 
    \item aim to make MLR.4 plausible
    \item \textbf{not} be based primarily on statistical tools ($R^2$, adjusted $R^2$, $t$-tests, $F$-tests.)
\end{itemize}
\bigskip \color{gray}
If our aim is prediction rather than causal analysis, statistical tools can be very helpful. \\
Aim: explain $Y$ as good as possible  
\end{frame}

\begin{frame}{Choice of regressors in causal analysis}
     
We are typically interested in the causal effect of one specific variable $X_1$. \\
What other variables $X_j$, $j = 2,…,k$ should we add as control variables?

\begin{itemize}
    \item Good control variables: \color{teal}\textbf{help to avoid OVB} \color{black}
    \item Bad control variables: \color{orange}\textbf{mediator variables} \color{black}(could be considered as an outcome variable themselves)

\end{itemize}
\end{frame}
\begin{frame}{Example for bad control variable}
Suppose, we are interested in estimating the causal effect of agricultural pesticide use on health care expenditures. 

$Health \: Expenditures_r = \beta_0 + \beta_1 pesticides_r + … + u_i$, $r = 1,…,R$ (region)

\bigskip
Should we control for the number of doctor visits made? 

\pause

No!
An increase in the number of doctor visits is a \alert{\textbf{channel}} through which pesticide use increases health care expenditures.  

If we control for the number of doctor visits, $\beta_1$ would reflect the impact of pesticide use on health care expenditures that are not generated by more doctor visits $\rightarrow$ a parameter that is generally  not of pratical interest.
\end{frame}
\end{comment}

\section{Recap and Next Steps}

\begin{frame}{Recap}
\begin{itemize}
    \item We are interested in estimating \alert relationships between economic variables.

\item This requires that we interpret the linear model as some stylized (economic) model, where the $\beta$s describe causal effects of $\boldsymbol{X}$ on $Y$.

\item When we interpret coefficients in that way, the \alert{Exogeneity Assumption} is crucial to get consistent estimates of the causal effects with OLS.

\item  OLS estimators $\hat{\beta_j}$ consistently estimate the \textbf{(partial) correlation} between $X_j$ and $Y$. But if there is \textbf{endogeneity}, this  does not represent the underlying causal relationship.

\item Endogeneity can arise because of reverse causality, measurement error and omitted variables. It is an extremely frequent problem in the social sciences.
\end{itemize}
\end{frame}

\begin{frame}{Next Steps}
How can we address the endogeneity issue?
\begin{enumerate}
    \item Instrumental Variables.
    \item Panel Data.
    \item Differences in Differences.
    \item Regression Discontinuity Design.
\end{enumerate}
\end{frame}
\appendix
\begin{frame}

\Large \centering
APPENDIX
    
\end{frame}
\begin{frame}{What do slope coefficients in a MLR measure? \\
The Frisch-Waugh Theorem}

In a MLR, $\hat{\beta}_j$ reflects the partial correlation between $X_j$ and $Y$, after the effects of other regressors have been taken out. 
\begin{align*}
   Y_i =\beta _0 + \beta_1X_{1i} + \boldsymbol{X}'_{2i} \boldsymbol{\beta_2} + u_i , i = 1,…, n 
\end{align*}

\begin{enumerate}
    \item Regress $X_{1i}$ on  $\boldsymbol{X}_{2i} $: yields residuals $\hat{e}_{x_1i}$
    \item Regress $Y_i$ on  $\boldsymbol{X}_{2i}$: yields residuals $\hat{e}_{yi}$
    \item Regress $\hat{e}_{yi}$ on $\hat{e}_{xi}$
\end{enumerate}
$\rightarrow$ estimated slope parameter in step 3 is identical to the OLS estimator $\hat{\beta}_1$  in the MLR! (Frisch-Waugh Theorem)

\end{frame}

\begin{frame}{The Frisch-Waugh Theorem}

The error $\hat{e}_{x_1i}$ can be thought of as the component (or variation) of $X_{1i}$ that is not linearly explained by $\boldsymbol{X}_{2i} $.

Thus, the coefficient $\beta_1$ is a (rescaled) measure of linear correlation of $X_{1i}$ and $Y_i$, after \textbf{partialling out (controlling linearly for) the effects of all variables in}  $\boldsymbol{X}_{2i} $


$\rightarrow$ Basis for the \textbf{ceteris paribus interpretation} of the model coefficients in a multivariate regression set-up. 

But: we keep only the other included regressors fixed. Causal inference requires that there are no components left in the error that are systematically correlated with $X_1$.
\end{frame}
\end{document}
